\documentclass[letterpaper]{article}
\usepackage{aaai}

\usepackage{times}
\usepackage{helvet}
\usepackage{courier}

\usepackage{graphicx}

\pdfinfo{
/Title (YARP)
/Subject (YARP)
/Author (IIT)}
% The file aaai.sty is the style file for AAAI Press 
% proceedings, working notes, and technical reports.
%
\title{YARP, a Thin Middleware for (Humanoid) Robots}
\author{Paul Fitzpatrick \and Giorgio Metta \and Lorenzo Natale \\
Address}
\setcounter{secnumdepth}{0}

\begin{document} 
\maketitle
\begin{abstract}
\begin{quote}

YARP stands for ``Yet Another Robot Platform.''  It is a robot
middleware that began life in 2000 as a thin layer over the QNX
real-time operating system to adapt it for use by humanoid robots.  It
is now used on all kinds of robots and operating systems around the
world.  YARP's communication model lies at a sweet spot that combines
efficiency, flexibility, and ease of use.

\end{quote}
\end{abstract}

\noindent 

YARP is plumbing for robot software.  YARP supports building a robot
control system as a collection of programs communicating in a
peer-to-peer way, with a family of connection types (tcp, udp,
multicast, local, MPI, mjpg-over-http, XML/RPC, tcpros, ...) that can
be swapped in and out to match your needs.  We also support similarly
flexible interfacing with hardware devices. Our strategic goal is to
increase the longevity of robot software projects
\cite{fitzpatrick08towards}.

YARP was, in its early days, shaped by the problem of keeping research
going despite constant flux in our robot platforms (hence the name).
It remains a ``reluctant'' middleware, with no desire or expectation
to be in control of your system.  This is from long experience telling
us that our systems always change.  
We want to minimize problem
of incompatible "architectures", "frameworks", and "middleware" (also
known in this context as "muddleware").

\begin{figure}
\centerline{\includegraphics[width=5cm]{icub.jpg}}
\caption{YARP is used on all sorts of robots,
from complex humanoids such as the iCub (pictured here),
to embedded systems, to UAVs.
} 
\end{figure}

YARP is not an operating system for your robot. We figure you already
have an operating system, or perhaps several. Nor does it do package
management (we like the package managers we have). We're not out for
world domination. It is easy to interoperate with YARP-using programs
without yourself necessarily having to use YARP. 
YARP is written in C++, with a core that uses
no external libraries, not even the standard template libraries, with
the exception of a small portion of ACE for portability (and this
portion can easily be embedded). YARP is free and open.  The
core YARP libraries are released under the LGPLv2.1.


\section{A potted history of YARP}

YARP was born in 2000 on an early humanoid robot (called Kismet)
controlled by a set of Motorola 68332 processors, an Apple Mac, and a
loose network of PCs running QNX, Linux, and Microsoft Windows
\cite{metta06yarp}.  Communication on this robot was a cocktail
mixture of dual-port RAM polling, QNX message passing, CORBA, and raw
sockets.

YARP worked very well on QNX, and decently on Linux.  For Microsoft
Windows, Macs, and other platforms an important development in YARP's
history was to take on a dependency on the ACE library [ref].  This
happened around 2003, with the immediate goal of simplifying the
addition of multicast support.  ACE has proven its worth to us many
times over for cross-platform networking, although we needed
to carefully keep any reference to ACE out of public header files
in order to avoid inheriting some of its less desirable properties
(a somewhat unstable API, and constraints on header file inclusion
order).

YARP developers were early and enthusiastic adopters of CMake
\cite{fitzpatrick10cmaking}.  In 2006, we happily dumped
earlier custom build scripts in favor of CMake project descriptions,
and never looked back.  CMake was this missing piece for having a
truly comfortable cross-platform experience.

About the same time, we started using SWIG [ref] to support languages
other than C++.

Use of CMake and SWIG has become common in open source libraries.
The help users exploit software on many operating systems and 
languages.  By keeping YARP as ``just a library,'' there is no
obstacle to doing the same with YARP.  But since YARP facilitates
intercommunication, this has the powerful effect of allowing 
programs written in all these environments to communicate with
each other, and easily.

YARP goes further... see the ``TELNET'' test.

\section{Known carriers}

\begin{tabular}{|l|p{7cm}|}
\hline
\multicolumn{2}{|c|}{YARP carriers} \\
\hline
tcp & Regular tcp \\
fast\_tcp & Variant that drops flow control \\
udp & UDP \\
http & Basic web interface \\
mcast & Multicast - avoid repeating the same data going
to many clients ``on the wire''  \\
local & local \\
mpi & delegate to MPI (Daniel Krieg) \\
xmlrpc & translate messages into XML/RPC compatible form \\
tcpros & interoperate with ROS publisher/subscribers \\
mjpeg & receive/transmit images in mjpeg-over-http format \\
text & send messages in human readable plain text form \\
shmem & use shared memory \\
\hline
\end{tabular}

Similarly, many devices.

\section{Progressive Enhancement}

\section{Data liberation}

* Google has a subproject to make sure that it is easy for users to copy all their data from a service in usable form, and migrate to a non-Google service: the "data liberation front"

* For YARP, we try to make it easy to users to redirect data streams to non-YARP based programs.  This lowers the cost of "boundary problems" - hooking up an unsupported device, computer, network, ...

* The protocol used for a YARP connection is decided at connection-time.  The initiator of the connection is free to choose from a wide set of protocols, and so can optimize for simplicity (using a plain text protocol), speed (udp), scalability (multicast), etc.  

* A YARP protocol need not be specific to YARP.  For example, YARP supports mjpeg-over-http.  This means that a YARP image source can be viewed directly from a browser without bridging, or a YARP image sink can receive data directly from a jpeg-streaming IP camera without bridging.  YARP also supports XML/RPC, so can connect to certain websites (or act as a webserver).

* YARP connections can be initiated by either the "sender" or "receiver", since the logical flow of data can be freely reversed.  This is important for supporting a wide range of protocols, which may be "pull" or "push" in nature.  It also makes connecting to YARP programs without using the YARP libraries even easier.  With other middleware, you get stuck having to make a server for at least one of the directions of data flow.

* A YARP network is designed to be usable without YARP.  Examples of making YARP connections without using the YARP library are available in C and Python.


\section{The TELNET test}

The TELNET test is this: can a user monitor and insert traffic on a
YARP network using just a telnet client (``telnet ip.address
port\_number'')?  In other words, without any use of YARP libraries or
programs?  This test has evolved out of past frustration with other
middleware, where simply passing a few numbers to a collaborator's
program can require jumping through a dozen hoops.  We try to follow
the ``golden rule,'' by making it really easy for others to send a few
numbers to a YARP-based program without having to dig through protocol
specifications, or link against our libraries, or use our build
machinery~-- all of which can be much more expensive in developer time
than one might think, given the variety of languages, operating
systems, development environments, and versions-of-everything in use.

Some keys to acing the TELNET test:

\begin{itemize}

\item TCP connections should be supported.  
  This is true of most middleware operating on regular
  networks. Some just support UDP, so this test doesn't
  make sense for them.

\item The basic communication model shouldn't stray too far from the
  notion of making ``connections'' to a named destination.

\item The direction in which a connection is initiated shouldn't
  determine the direction of data flow.

\item Human readable/writable.

\end{itemize}


\section{YARP plus/minus}

YARP core libraries and their full dependencies are light.
YARP has for example been used on an embedded system
with 54MB RAM, 28MB Flash.

YARP is portable, serving an interdisciplinary community.

No code generation, no special build system, doesn't try to take over.

Flexible, open model of connections - proving good for evolution and interoperation.

Passes the ``telnet test''.

Negatives:

Native implementation exists only in C++.  

Diverse OS support limits growth of library, since any
new dependency is very costly.

Lack of IDL / code generation can lead to some tedium
implementing classic RPC-style code.


\bibliography{yarp}
\bibliographystyle{aaai}


\end{document}
